% ===== PRÉAMBULE CANONIQUE — à coller VERBATIM en tête de CHAQUE .tex =====
\documentclass[11pt,a4paper]{article}
\usepackage[utf8]{inputenc}\usepackage[T1]{fontenc}\usepackage{lmodern}
\usepackage[french]{babel}
\usepackage{amsmath,amssymb,mathtools}\usepackage{icomma}
\usepackage{siunitx}\sisetup{locale=FR}  % \qty{}{}, \si{}, \ang{}
\usepackage{xcolor}\usepackage{colortbl}
\usepackage{tikz}\usepackage{pgfplots}\pgfplotsset{compat=1.18}
\usepackage[siunitx,europeanresistors]{circuitikz} % circuits (élec.)
\usepackage[most]{tcolorbox}\usepackage{enumitem}\usepackage{array,booktabs}
\usepackage[a4paper,margin=2.2cm]{geometry}
\usetikzlibrary{arrows.meta,calc,angles,quotes,positioning,patterns,%
  backgrounds,shapes.geometric,decorations.pathmorphing,decorations.markings,intersections}
\definecolor{primary}{RGB}{222,49,99}\definecolor{primarydark}{RGB}{150,22,55}
\definecolor{defcol}{RGB}{0,114,178}\definecolor{methcol}{RGB}{52,92,148}
\definecolor{excol}{RGB}{0,135,90}\definecolor{attncol}{RGB}{200,40,55}
\tcbset{boxbase/.style={enhanced,breakable,boxrule=0.9pt,arc=2.5pt,
  left=3mm,right=3mm,top=2.3mm,bottom=2.3mm,fonttitle=\bfseries,coltitle=white}}
% --- boîtes TUTORIEL (noms EXACTS, ne pas renommer) ---
\newtcolorbox{cle}[1][]{boxbase,colback=defcol!6,colframe=defcol,
  title={\textsc{À retenir}\ifx\relax#1\relax\relax\else\textmd{ : \itshape #1}\fi}}
\newtcolorbox{methode}[1][]{boxbase,colback=methcol!7,colframe=methcol,sharp corners,
  title={\textsc{Méthode}\ifx\relax#1\relax\relax\else\textmd{ --- \itshape #1}\fi}}
\newtcolorbox{exemplebox}[1][]{boxbase,colback=excol!5,colframe=excol,
  title={\textsc{Exemple}\ifx\relax#1\relax\relax\else\textmd{ : \itshape #1}\fi}}
\newtcolorbox{piege}[1][]{boxbase,colback=attncol!6,colframe=attncol,
  title={\textsc{Piège}\ifx\relax#1\relax\relax\else\textmd{ : \itshape #1}\fi}}
\newtcolorbox{remarque}[1][]{boxbase,colback=black!4,colframe=black!45,
  title={\textsc{Remarque}\ifx\relax#1\relax\relax\else\textmd{ : \itshape #1}\fi}}
% --- boîtes EXERCICE / CORRIGÉ (noms EXACTS, auto-comptées) ---
\newtcolorbox[auto counter]{exo}[1][]{boxbase,colback=primary!4,colframe=primary,
  title={\textsc{Exercice}~\thetcbcounter\ifx\relax#1\relax\relax\else\textmd{ --- \itshape #1}\fi}}
\newtcolorbox[auto counter]{corrige}[1][]{boxbase,colback=excol!4,colframe=excol,
  title={\textsc{Corrigé}~\thetcbcounter\ifx\relax#1\relax\relax\else\textmd{ --- \itshape #1}\fi}}
\newcommand{\vv}[1]{\vec{#1}}\newcommand{\ds}{\displaystyle}
\newcommand{\R}{\mathbb{R}}\newcommand{\N}{\mathbb{N}}\newcommand{\Z}{\mathbb{Z}}
% (un \usepackage{fancyhdr}+en-tête est possible ; il est ignoré côté web)
% =========================================================================
\begin{document}
% THEME_TITLE: Le poids et l'équilibre de translation
\begin{center}
{\LARGE\bfseries\color{primarydark}Thème 1 — Le poids et l'équilibre de translation}\\[2pt]
{\color{primary}\itshape Statique — Fiche 3}
\end{center}
\vspace{1mm}

Une \textbf{force} $\vv F$ est une grandeur \emph{vectorielle} (point d'application, direction, sens,
\textbf{intensité} $F=\lVert\vv F\rVert$ en \si{\newton}). Un solide est en \textbf{équilibre de
translation} lorsque toutes les forces qu'il subit se compensent : il ne se déplace pas.

\begin{cle}[le poids $\vv G = m\,\vv g$]
Le \textbf{poids} est la force d'attraction qu'exerce la Terre sur un corps :
$\boxed{\,G = m\,g\,}$ ($G$ en \si{\newton}, masse $m$ en \si{\kilogram},
$g\approx\qty{9.81}{N/kg}$ souvent arrondi à $\qty{10}{N/kg}$). La \emph{masse} $m$ est la même
partout ; le \emph{poids} dépend de l'astre — sur la Lune, $g$ est six fois plus faible, donc le poids
est six fois plus \textbf{petit} (jamais plus grand !), la masse restant inchangée.
\end{cle}

\begin{piege}[toujours convertir la masse en kilogrammes]
$G=mg$ n'est valable que si $m$ est en \si{\kilogram}. Pour $m=\qty{200}{g}=\qty{0.200}{kg}$ et
$g=\qty{10}{N/kg}$ : $G = 0{,}200\times 10 = \qty{2}{N}$ (et non \qty{200}{N} !). Oublier la
conversion fait une erreur d'un facteur \textbf{1000}.
\end{piege}

\begin{cle}[équilibre de translation : $\sum\vv F=\vv 0$]
Un solide soumis aux forces $\vv F_1,\dots,\vv F_n$ est en équilibre de translation si la
\textbf{résultante} (somme \emph{vectorielle}) est nulle : $\boxed{\,\sum_i \vv F_i = \vv 0\,}$.
Ce n'est pas la somme des \emph{intensités} qui s'annule (une intensité est positive) : les forces se
\emph{compensent}. Dans un plan : $\sum F_x = 0$ \textbf{et} $\sum F_y = 0$. C'est aussi la condition
d'un mouvement à \textbf{vitesse constante} ($\vv a=\vv 0$).
\end{cle}

\begin{exemplebox}[réaction d'un support]
Un livre immobile sur une table subit son poids $\vv G$ (bas) et la \textbf{réaction normale} $\vv R$
(haut). L'équilibre $\vv G+\vv R=\vv 0$ donne $R=G$, même direction, sens opposé.
\hfill
\begin{tikzpicture}[>=Stealth,scale=0.75,baseline=-2mm]
  \draw[black!55,line width=0.9pt] (-1.5,0)--(1.5,0);
  \foreach \x in {-1.3,-0.9,...,1.4}\draw[black!45](\x,0)--(\x-0.2,-0.25);
  \draw[fill=primary!18,draw=primary!70] (-0.55,0) rectangle (0.55,0.45);
  \fill (0,0.22) circle (1.1pt);
  \draw[->,line width=1pt,defcol] (0,0.22)--(0,1.4) node[above]{$\vv R$};
  \draw[->,line width=1pt,orange] (0,0.22)--(0,-0.95) node[below]{$\vv G$};
\end{tikzpicture}
\end{exemplebox}

\begin{cle}[objet suspendu à des fils]
Un objet suspendu immobile a son poids \emph{intégralement} porté par les fils : la somme des
\textbf{composantes verticales} des tensions vaut $G$. Pour \textbf{deux fils symétriques}, chacun en
porte la moitié, $T=G/2$. En montage \emph{asymétrique}, le fil le plus proche de la verticale
supporte davantage — mais la \emph{somme} reste égale à $G$.
\end{cle}

\begin{piege}[action/réaction $\neq$ équilibre]
3\textsuperscript{e} loi de Newton : si $A$ exerce $\vv F$ sur $B$, alors $B$ exerce $\vv F'=-\vv F$ sur
$A$. Ces forces opposées \emph{s'appliquent sur deux corps différents} : elles ne s'équilibrent jamais
entre elles. Pour l'équilibre d'\emph{un} corps, on ne retient que les forces exercées \emph{sur lui}.
\end{piege}

\end{document}
